% !TEX root =./ECE444_Practice_main.tex
%
% L24 practice set (Sidelobes and Tapering Theory) -- why an untapered aperture
% has sidelobes, the cosine/Chebyshev/Taylor families, the three costs of a
% taper, and converting a design to PHASER element gains.

\fullwidth{\textbf{Learning Objective 3.7: } I can apply amplitude tapering
(uniform, cosine, Chebyshev, Taylor) to control sidelobe level and predict the
pattern trade-off.
\\[4pt]\normalfont\small Show your work. A \textbf{Documentation} line at the end
is required (write \textbf{None} if you did not collaborate).}

\begin{questions}

\question \textbf{Where sidelobes come from, and why a system cares.}
The PHASER's eight elements are driven at equal amplitude and spaced
$d = 14\text{ mm}$, which is $d/\lambda = 0.481$ at $10.3\ \ghz$. Its measured
first sidelobe sits $12.8\db$ below the peak.
\begin{parts}
	\part In two or three sentences, explain why a uniformly illuminated aperture
	has sidelobes at all. Use the Fourier relationship between the aperture
	illumination and the far-field pattern in your answer.
		\begin{solution}[1.4in]
		The far-field pattern is the Fourier transform of the aperture
		illumination, with the transform variable being the space frequency
		$k\sinp{\theta}$. Uniform illumination is a rectangle: full amplitude
		across the array and nothing beyond the ends. That abrupt edge is a
		discontinuity, and a discontinuity contains energy at every spatial
		frequency, so the transform cannot fall off quickly. The transform of a
		rectangle is a sinc function, whose ripples decay only as $1/u$, and
		those ripples are the sidelobes.
		\end{solution}

	\begin{minipage}[t]{0.58\linewidth}
		\part The array is rebuilt with $N = 16$ elements at the same spacing,
		still uniformly driven. Compute the new half-power beamwidth from
		$\theta_\text{HP} \approx 0.886\lambda/(Nd)$, and state what happens to the
		first sidelobe level.
		\begin{solution}[1.3in]
		\eq{ \theta_\text{HP} &\approx \frac{0.886}{16 \times 0.481} = 0.1151
		\text{ rad} = 6.6\mydeg }
		The first sidelobe does not move: it stays near $-13\db$. Doubling the
		aperture halves the beamwidth, but sidelobe level is set by the
		\emph{shape} of the illumination, not by its size, and the shape is still
		uniform.
		\end{solution}
	\end{minipage}\hfill%
	\begin{minipage}[t]{0.37\linewidth}
		\raggedleft \vspace{12pt}
		$\theta_\text{HP}$ = \ansbox[1.3in]{$6.6\mydeg$}
	\end{minipage}

	\begin{minipage}[t]{0.58\linewidth}
		\part A search radar looks at a target on boresight while ground clutter
		$40\db$ stronger than the target return arrives through the $-12.8\db$
		first sidelobe. Compute the clutter-to-target ratio at the beamformer
		output, then find the sidelobe level that would put the clutter $3\db$
		\emph{below} the target.
		\begin{solution}[1.4in]
		The target enters on the main beam at $0\db$ of pattern gain and the
		clutter enters $12.8\db$ down, so the antenna gives the target a
		$12.8\db$ advantage:
		\eq{ C/T &= 40 - 12.8 = 27.2\db \text{ (clutter above target)} }
		To land at $C/T = -3\db$ the pattern must supply $43\db$ of
		discrimination:
		\eq{ \text{SLL} &= -(40 + 3) = -43\db }
		\end{solution}
	\end{minipage}\hfill%
	\begin{minipage}[t]{0.37\linewidth}
		\raggedleft \vspace{12pt}
		$C/T$ = \ansbox[1.2in]{$27.2\db$}
		\\[10pt]
		SLL = \ansbox[1.2in]{$-43\db$}
	\end{minipage}
\end{parts}

\newpage

\question \textbf{The Hann preset by the numbers.}
The PHASER GUI's Hann preset sets the eight Element Gains to
$12, 43, 77, 100, 100, 77, 43, 12$ percent. Work in normalized amplitudes
$a_n$ of $0.12$, $0.43$, $0.77$, $1.00$, $1.00$, $0.77$, $0.43$, $0.12$.
\begin{parts}
	\begin{minipage}[t]{0.58\linewidth}
		\part Compute $\sum a_n$ and $\sum a_n^2$.
		\begin{solution}[1.3in]
		The taper is symmetric, so sum the inner four and double:
		\eq{ \sum a_n &= 2\lp 0.12 + 0.43 + 0.77 + 1.00 \rp = 2\lp 2.32 \rp = 4.64 \\
		\sum a_n^2 &= 2\lp 0.0144 + 0.1849 + 0.5929 + 1.0000 \rp = 3.584 }
		\end{solution}
	\end{minipage}\hfill%
	\begin{minipage}[t]{0.37\linewidth}
		\raggedleft \vspace{12pt}
		$\sum a_n$ = \ansbox[1.1in]{$4.64$}
		\\[10pt]
		$\sum a_n^2$ = \ansbox[1.1in]{$3.584$}
	\end{minipage}

	\begin{minipage}[t]{0.58\linewidth}
		\part Compute the taper efficiency $\eta_t = \lp\sum a_n\rp^2 / \lp N \sum a_n^2 \rp$
		as a ratio and in decibels.
		\begin{solution}[1.3in]
		\eq{ \eta_t &= \frac{\lp 4.64 \rp^2}{8 \times 3.584} = \frac{21.53}{28.67}
		= 0.751 \\
		10\mylog{0.751} &= -1.24\db }
		\end{solution}
	\end{minipage}\hfill%
	\begin{minipage}[t]{0.37\linewidth}
		\raggedleft \vspace{12pt}
		$\eta_t$ = \ansbox[1.5in]{$0.751 = -1.2\db$}
	\end{minipage}

	\begin{minipage}[t]{0.58\linewidth}
		\part Compute how far the peak of the measured sweep falls below the
		uniform-taper trace.
		\begin{solution}[1.2in]
		\eq{ \text{peak drop} &= 20\mylog{\frac{\sum a_n}{N}}
		= 20\mylog{\frac{4.64}{8}} \\
		&= 20\mylog{0.580} = -4.7\db }
		\end{solution}
	\end{minipage}\hfill%
	\begin{minipage}[t]{0.37\linewidth}
		\raggedleft \vspace{12pt}
		peak drop = \ansbox[1.3in]{$-4.7\db$}
	\end{minipage}

	\part Parts (b) and (c) are both decibel losses, and they are not equal. In
	two or three sentences, say what each one measures and which of the two is
	the loss of directivity.
		\begin{solution}[1.5in]
		The taper efficiency in part (b) is the fraction of the uniform array's
		directivity that the tapered array keeps, so $-1.2\db$ is the loss of
		antenna performance. The peak drop in part (c) is a coherent
		receive-voltage loss: six of the eight channels were attenuated, so the
		plotted peak fell, but the channel noise was attenuated along with the
		signal. The difference,
		$-4.7 - \lp -1.2 \rp = -3.5\db$, is gain that was simply turned down and
		would return if the sliders could be scaled back up past $100\%$.
		\end{solution}
\end{parts}

\newpage

\question \textbf{A Chebyshev design to specification.}
A customer specifies that no sidelobe may exceed $-25\db$ relative to the peak,
on the eight-element array at $d/\lambda = 0.481$. The uniform reference is
$\theta_\text{HP} = 13.2\mydeg$.
\begin{parts}
	\begin{minipage}[t]{0.58\linewidth}
		\part Compute the voltage ratio $R$ and the Chebyshev stretch factor
		$x_0 = \cosh\lp \cosh^{-1}(R)/(N-1) \rp$.
		\begin{solution}[1.4in]
		\eq{ R &= 10^{25/20} = 10^{1.25} = 17.78 \\
		\cosh^{-1}(17.78) &= 3.571 \\
		x_0 &= \cosh\lp \frac{3.571}{7} \rp = \cosh\lp 0.5101 \rp \\
		&= 1.133 }
		\end{solution}
	\end{minipage}\hfill%
	\begin{minipage}[t]{0.37\linewidth}
		\raggedleft \vspace{12pt}
		$R$ = \ansbox[1.1in]{$17.78$}
		\\[10pt]
		$x_0$ = \ansbox[1.1in]{$1.133$}
	\end{minipage}

	The Dolph expansion for this $x_0$ returns element amplitudes of
	$38, 58, 84, 100, 100, 84, 58, 38$ percent.

	\begin{minipage}[t]{0.58\linewidth}
		\part Compute the taper efficiency and the peak drop for this design.
		\begin{solution}[1.5in]
		\eq{ \sum a_n &= 2\lp 0.38 + 0.58 + 0.84 + 1.00 \rp = 5.60 \\
		\sum a_n^2 &= 2\lp 0.1444 + 0.3364 + 0.7056 + 1 \rp \\
		&= 4.373 \\
		\eta_t &= \frac{\lp 5.60 \rp^2}{8 \times 4.373} = \frac{31.36}{34.98}
		= 0.896 = -0.48\db \\
		\text{peak drop} &= 20\mylog{\frac{5.60}{8}} \\
		&= 20\mylog{0.700} = -3.1\db }
		\end{solution}
	\end{minipage}\hfill%
	\begin{minipage}[t]{0.37\linewidth}
		\raggedleft \vspace{12pt}
		$\eta_t$ = \ansbox[1.4in]{$0.896 = -0.48\db$}
		\\[10pt]
		peak drop = \ansbox[1.2in]{$-3.1\db$}
	\end{minipage}

	\begin{minipage}[t]{0.58\linewidth}
		\part This design has a beam broadening factor of $1.20$. Compute its
		half-power beamwidth.
		\begin{solution}[1.1in]
		\eq{ \theta_\text{HP} &= 1.20 \times 13.2\mydeg = 15.8\mydeg }
		\end{solution}
	\end{minipage}\hfill%
	\begin{minipage}[t]{0.37\linewidth}
		\raggedleft \vspace{12pt}
		$\theta_\text{HP}$ = \ansbox[1.2in]{$15.8\mydeg$}
	\end{minipage}

	\part State the price of this design for the customer in one sentence, then
	explain in one or two more why a Chebyshev design beats a Hann window that
	reaches a similar sidelobe level.
		\begin{solution}[1.6in]
		Sidelobes drop from $-12.8\db$ to $-25\db$, and the cost is a beam
		$2.6\mydeg$ wider, $0.48\db$ of directivity, and a $3.1\db$ lower peak on
		the displayed trace. The Chebyshev design is optimal in the sense that no
		other set of amplitudes reaches $-25\db$ with a narrower main lobe,
		because it holds every sidelobe exactly at the specification. A Hann
		window pushes its far sidelobes well below $-25\db$, which is suppression
		the specification never asked for, and it paid for that suppression in
		beamwidth it did not need to spend.
		\end{solution}
\end{parts}

\newpage

\question \textbf{Reading a tapered sweep.}
A student loads an unknown taper into the GUI and runs a beam sweep against the
HB100. Compared with the uniform trace taken a minute earlier, the peak is
$3.8\db$ lower, the half-power beamwidth reads $17.5\mydeg$, and no sidelobe is
visible above the sweep's noise floor, which sits $23\db$ below the
\emph{uniform} peak.
\begin{parts}
	\part Which of the designs from this lesson is the student most likely
	using? Give the two measurements that identify it.
		\begin{solution}[1.3in]
		The $-30\db$ Chebyshev design, whose element gains are
		$26, 52, 81, 100, 100, 81, 52, 26$ percent. It predicts a peak drop of
		$-3.8\db$ and a half-power beamwidth of $17.1\mydeg$, and the measured
		$-3.8\db$ and $17.5\mydeg$ match both within the sweep's angular
		resolution. The Hann preset would have dropped the peak $4.7\db$ and
		widened the beam to about $19.5\mydeg$.
		\end{solution}

	\begin{minipage}[t]{0.58\linewidth}
		\part How much directivity did the array lose? The design has
		$\sum a_n = 5.18$ and $\sum a_n^2 = 3.988$.
		\begin{solution}[1.3in]
		\eq{ \eta_t &= \frac{\lp 5.18 \rp^2}{8 \times 3.988} = \frac{26.83}{31.91}
		= 0.841 \\
		10\mylog{0.841} &= -0.75\db }
		\end{solution}
	\end{minipage}\hfill%
	\begin{minipage}[t]{0.37\linewidth}
		\raggedleft \vspace{12pt}
		loss = \ansbox[1.2in]{$-0.75\db$}
	\end{minipage}

	\part The student writes in the lab notebook that the array lost $3.8\db$ of
	gain. Correct that statement in two sentences.
		\begin{solution}[1.4in]
		The array lost $0.75\db$ of directivity, which is the taper efficiency
		computed in part (b). The remaining $3.0\db$ of the drop is signal that
		six attenuated channels never delivered to the summing node, and the
		noise from those channels was attenuated by the same factor, so the
		sensitivity of the array did not fall by $3.8\db$.
		\end{solution}

	\begin{minipage}[t]{0.58\linewidth}
		\part State the tightest upper bound the student can honestly put on this
		taper's sidelobe level, expressed relative to the taper's own peak.
		\begin{solution}[1.3in]
		The floor sits $23\db$ below the uniform peak, and this taper's peak is
		itself $3.8\db$ below the uniform peak, so relative to its own peak the
		floor is only
		\eq{ 23 - 3.8 &= 19.2\db \text{ down} }
		The sidelobes are therefore below $-19\db$ relative to the tapered peak.
		The theoretical value is $-30\db$, but the measurement cannot confirm
		anything deeper than the floor allows.
		\end{solution}
	\end{minipage}\hfill%
	\begin{minipage}[t]{0.37\linewidth}
		\raggedleft \vspace{12pt}
		SLL $<$ \ansbox[1.3in]{$-19.2\db$}
	\end{minipage}
\end{parts}

\newpage

\question \textbf{Design judgment on eight elements.}
A program office asks for $-50\db$ sidelobes from the eight-element PHASER
array. The Chebyshev design that meets it has element amplitudes
$9$, $35$, $72$, $100$, $100$, $72$, $35$, $9$ percent, a beam broadening factor of $1.51$,
and $\eta_t = 0.710$.
\begin{parts}
	\begin{minipage}[t]{0.58\linewidth}
		\part Compute the half-power beamwidth and the peak drop for this design,
		using $13.2\mydeg$ as the uniform reference.
		\begin{solution}[1.4in]
		\eq{ \sum a_n &= 2\lp 0.09 + 0.35 + 0.72 + 1.00 \rp = 4.32 \\
		\theta_\text{HP} &= 1.51 \times 13.2\mydeg = 19.9\mydeg \\
		\text{peak drop} &= 20\mylog{\frac{4.32}{8}} \\
		&= 20\mylog{0.540} = -5.3\db }
		\end{solution}
	\end{minipage}\hfill%
	\begin{minipage}[t]{0.37\linewidth}
		\raggedleft \vspace{12pt}
		$\theta_\text{HP}$ = \ansbox[1.2in]{$19.9\mydeg$}
		\\[10pt]
		peak drop = \ansbox[1.2in]{$-5.3\db$}
	\end{minipage}

	\part The ADAR1000 sets each element gain to a resolution of about one
	percent of full scale. Explain in two or three sentences why that resolution
	threatens a $-50\db$ design more than it threatens a $-25\db$ design.
		\begin{solution}[1.6in]
		A sidelobe at $-50\db$ is a residue about three parts in a thousand of
		the peak, so it survives only if the eight amplitudes hold their designed
		ratios to better than that. The outermost elements of this design sit at
		$9\%$, where a one-percent step is an eleven-percent error on that
		element, which is two orders of magnitude larger than the residue the
		design is trying to preserve. The $-25\db$ design keeps its smallest
		element at $38\%$, where the same step is under three percent, so its
		sidelobes move by a fraction of a decibel instead of being swamped.
		\end{solution}

	\part Recommend what to do about the program office's request. Your answer
	must name the specification you would propose instead and give the reason in
	terms of what the hardware can show.
		\begin{solution}[1.7in]
		Propose $-30\db$ rather than $-50\db$. The beam sweep's noise floor is
		about $23\db$ below the uniform peak, so a $-30\db$ design is already
		below anything this instrument can measure, and a $-50\db$ design cannot
		be verified on the bench at all. The deeper specification also costs
		$3\mydeg$ more beamwidth and roughly $0.7\db$ more directivity than the
		$-30\db$ design, and its edge elements are too small to set accurately.
		If the office needs suppression in one specific direction rather than
		everywhere, the correct tool is a steered null, not a deeper taper.
		\end{solution}
\end{parts}

\vspace{1.5em}
\noindent\textbf{Documentation:} \rule[-2pt]{0.6\linewidth}{0.4pt}

\end{questions}
